\section{Introduction}\label{sec:intro}

Activation steering modifies language-model behavior by adding a contrastively constructed vector to the residual stream at a middle layer. This intervention can shift benchmark measures of honesty, refusal, or sentiment in the intended direction \citep{zou2023repe,rimsky2024caa,li2023iti,arditi2024refusal}. Because it targets an intermediate representation, the shift is often interpreted as a change to a high-level concept used by the model. For honesty, that interpretation supports a safety claim: a vector that changes an internal representation of truthfulness should generalize across phrasings, domains, and languages.

The same benchmark gain has a simpler explanation. A vector that shifts the output distribution toward words such as \emph{true} and \emph{honest}, and away from \emph{lie} and \emph{deceive}, could improve the score without altering upstream computation. An upstream concept change and a readout-level vocabulary bias can produce the same TruthfulQA score but should behave differently out of distribution. Standard evaluation cannot separate them because an intervention at layer $\ell^{*}$ changes a shared residual consumed by both the unembedding and the remaining transformer blocks. Confusing the two mechanisms turns evidence about word probabilities into a claim about internal computation and overstates what a benchmark gain establishes.

We separate the two paths directly. The first-order logit response to an injected vector $v$ consists of a direct term $\WUt^{*} v$, consumed unchanged by the readout, and an indirect term $\WUt^{*} M v$, propagated through downstream attention and MLP blocks (Section~\ref{sec:method}). For a selected token set $T$, we project the steering vector onto the orthogonal complement of the corresponding rows of the \emph{effective} unembedding, which composes the unembedding with the RMSNorm Jacobian at the readout point. Under the calibration-averaged effective unembedding, the projected vector $\vperp$ has zero first-order direct effect on every token in $T$, certified to machine precision. We use the behavioral fraction that survives this excision, $\rho = \Delta(\vperp)/\Delta(\vdec)$, as a depth statistic: values near $0$ indicate direct readout through $T$, while values near $1$ indicate downstream computation under this decomposition.

The depth statistic is informative only when two conditions hold. Its denominator must be a genuine steering effect, so we select each operating point under a coherence gate that rejects magnitudes with excessive held-out perplexity. Without this gate, a large coefficient can raise teacher-forced multiple-choice scores while reducing free generation to noise, making $\rho$ a ratio of two noise terms. We also compare aligned excision with equal-rank random removal. This contrast determines whether any loss is specific to the aligned readout or reflects generic projection damage.

On Llama-3-8B-Instruct, contrastive activation addition (CAA) strongly shifts honesty-coded vocabulary, but its validation-selected coherent point yields no statistically significant positive truthfulness gain; the largest gain elsewhere in the tested coherent grid is less than one fifth of mass-mean's. The mass-mean direction improves truthfulness (test-set $\Delta$MC2 of \MmTestDelta{} \MmTestDeltaCI), and the gain survives certified aligned-64 readout excision ($\rho = \MmRhoAlSixtyFour$ \MmRhoAlSixtyFourCI). In-domain, the loss is indistinguishable from random-subspace controls; the projected gain also survives paraphrase shift and free-form generation. At the selected operating points across four models, CAA yields no statistically significant positive gain. Mass-mean succeeds on Llama-3, Mistral, and Llama-2 but not Qwen, and most of each successful gain survives readout excision. Useful steering is model-dependent, but successful instances are predominantly downstream under the tested readout rather than vocabulary suppression.

\paragraph{Contributions.}
We make three contributions. (i) We introduce a token-conditional projection that zeroes a steering vector's calibration-averaged first-order direct logit contribution on a chosen token set, includes the RMSNorm correction omitted by prior unembedding-orthogonal constructions, and provides machine-precision certificates. (ii) We make the resulting depth statistic $\rho$ interpretable with a coherence-gated operating-point protocol and a random-projection control. (iii) We evaluate CAA and mass-mean across four models, separating whether a gain exists from whether it depends on the tested direct readout. Code, token lists, certificates, and the formal proof are available at \url{https://github.com/aryan-cs/liar-liar}.
